\documentclass{situationist-preprint}

\title{Duties to the Dead}
\msid{EXI-13-a}
\msfamily{EXI-13}
\msdiscipline{Meaning, Death and the Absurd}
\mstype{Principal manuscript}
\mskeywords{duties to the dead; posthumous harm; surviving interests; wills; reliance; promises; Feinberg; Pitcher; Partridge}

\begin{document}
\maketitle

\begin{abstract}
What do we owe the dead? We argue that we owe them fidelity to their reasonable wishes, and that this duty weakens over time. The account has two parts. The first concerns the subject of the duty. Following Feinberg and Pitcher, we take the duty to be owed to the person as she was before death, whose interests in what happens after her death are thwarted or fulfilled by posthumous events; on the tenseless view of time this involves no backward causation. The second concerns its ground. We extend the reliance account of promissory obligation defended in ETH-16: a person who lives in a community that honors the wishes of its dead is entitled to rely on its doing so, and plans her life on that reliance; the duty to the dead is owed to that entitlement, created while she lived. The entitlement protects interests in projects, dependents and reputation, and these interests are progressively completed, extinguished or overtaken as time passes. The duty does not lapse, but its weight against the claims of the living diminishes with the interests it protects. The account explains why only reasonable wishes bind, why deathbed promises bind more, and why the dead hand loses its grip.
\end{abstract}

\section{The problem}

We keep the promises we made to the dying, carry out wills, bury and cremate according to instructions, and think it wrong to defame the dead or to publish what they asked to be destroyed. These practices presuppose that something is owed to the dead. But the dead are not there to be owed anything. If, with Epicurus, nothing that happens after death can be bad for the one who has died, it is hard to see how a duty to the dead could be grounded in their good, and if duties of fidelity are owed to persons, it is unclear to whom such a duty is owed.

Three positions dominate the literature. \citet[ch.~2]{Feinberg1984} argued that a person's interests can survive her death, since many of them, in the success of her projects, the welfare of her children, the fate of her reputation, have objects that lie outside her experience and outside her lifetime; a posthumous event that thwarts such an interest harms her. \citet{Pitcher1984} identified the subject of posthumous harm as the ante-mortem person, the person as she was while alive, and argued that the harm involves no backward causation: the posthumous event makes it true that the ante-mortem person was harmed, without changing anything in the past. \citet{Partridge1981} denied that the dead have interests, and argued that our practices of respect for the dead are justified by the interests of the living, who need assurance that their own wishes will be honored. \citet{Callahan1987} argued that the dead cannot be harmed and that the appearance of posthumous harm is an illusion produced by our concern for the persons we remember. The thought that what happens after death can affect how a life went is older: Aristotle asks whether the fortunes of a person's descendants can affect whether she was happy, and answers that they can, though only slightly \citep[I.10--11]{AristotleNE}.

We take from Feinberg and Pitcher the subject of the duty and from Partridge the role of assurance, and combine them with the account of fidelity in ETH-16. The result explains both the content of the duty and its weakening.

\section{Framework}

\begin{definition}[Reasonable wish]\label{d:wish}
A \emph{reasonable wish} of a person about what happens after her death is a wish that (i) she expressed, or that those who survive her are in a position to know she had; (ii) concerns a matter in which she had a legitimate interest, such as her body, her property, her work, her dependents or her reputation; and (iii) can be honored without wronging others or imposing on them burdens disproportionate to the interest it serves.
\end{definition}

\begin{principle}[Surviving interests]\label{p:surv}
A person has interests whose objects lie after her death. If a posthumous event thwarts such an interest, it is bad for the ante-mortem person, in the sense of the deprivation account of EXI-02: her life as a whole goes worse than it would otherwise have gone.
\end{principle}

\begin{principle}[Reliance]\label{p:rel}
If $A$ voluntarily creates in $B$ a reasonable reliance on $A$'s doing $X$, knowing that $B$ wants assurance that $X$ will be done, then $A$ is obligated to $B$ to do $X$, unless $B$ releases $A$ or the obligation is defeated by weightier reasons. Such an obligation persists until it is discharged, released or defeated (ETH-16).
\end{principle}

\begin{assumption}[Practice]\label{a:practice}
In communities in which the wishes of the dead are honored, those who maintain the practice, by executing wills, keeping deathbed promises and enforcing testamentary law, voluntarily give each member a reasonable reliance that her reasonable wishes will be honored after her death, knowing that members want this assurance; and members plan their lives on it.
\end{assumption}

\begin{assumption}[Weight]\label{a:weight}
The weight of an obligation owed to a person's entitlement, against competing reasons, is proportional to the weight of the interests of that person which the entitlement protects.
\end{assumption}

Principle~\ref{p:surv} is Feinberg's thesis in Pitcher's form. The deprivation account of EXI-02 shows that an event can be bad for a person without her experiencing it and without her existing when it occurs (Lemma~3.2 of EXI-02); Principle~\ref{p:surv} applies the same structure to events that are bad for a person because they thwart her interests rather than because they cut her life short. On the tenseless view of time defended in TSM-01, the fact that a posthumous event thwarts an ante-mortem interest is a fact about the relation between the two, which obtains without being located at either time, and no event changes the past. Principle~\ref{p:rel} is the main result of ETH-16, which leaves open whether the dead can have entitlements. Assumption~\ref{a:practice} answers that question for the case that matters here: the entitlement is created while the person lives, and she relies on it while she lives, by writing a will rather than giving her property away, by entrusting her manuscripts to an executor, by making plans for her children that will be completed after she is gone.

\section{The argument}

\begin{lemma}[Subject]\label{l:subject}
Duties concerning the wishes of the dead are owed to the ante-mortem person.
\end{lemma}

\begin{proof}
By Assumption~\ref{a:practice}, the reliance on which the duty rests was created in the ante-mortem person, while she lived. By Principle~\ref{p:rel}, an obligation grounded in the creation of reasonable reliance is owed to the person in whom the reliance was created. And the interests the entitlement protects are, by Principle~\ref{p:surv}, interests of the ante-mortem person. So the duty is owed to her. That she does not exist when the duty is to be discharged is no obstacle, for the same reason that it is no obstacle to her being deprived by her death (EXI-02).
\end{proof}

\begin{lemma}[Content]\label{l:content}
The duty is one of fidelity to the person's reasonable wishes, and does not extend to unreasonable ones.
\end{lemma}

\begin{proof}
By Assumption~\ref{a:practice} and Principle~\ref{p:rel}, what the ante-mortem person is entitled to is that her reasonable wishes be honored. A wish that fails clause (iii) of Definition~\ref{d:wish} is one whose fulfilment would wrong others or burden them disproportionately; no practice can give a reasonable reliance on its fulfilment, since reliance on another's doing wrong is not reasonable. A wish that fails clause (i) cannot be relied on to be honored, since no one is in a position to honor it. A wish that fails clause (ii) concerns matters on which the person has no claim against the living. So the duty is fidelity to reasonable wishes.
\end{proof}

\begin{lemma}[Weakening]\label{l:weak}
As time passes after a person's death, the interests that her entitlement protects become, in the ordinary course, fewer and weaker.
\end{lemma}

\begin{proof}
The interests of Definition~\ref{d:wish}(ii) have objects that are themselves in time. Interests in the welfare of dependents are interests in the lives of particular persons, and they are fulfilled or extinguished as those persons grow up and die. Interests in the completion of projects are fulfilled when the projects are completed and lapse when their completion is no longer possible. Interests in reputation are interests in how one is regarded by those who knew one or knew of one, and their object narrows as those people die. Moreover, a wish is formed for the circumstances its holder could foresee, and as circumstances change, the wish's application to them is less and less determined by the interests that made it reasonable; a bequest to maintain a hospital for a disease that has been eradicated serves no interest its testator had in founding it. In each case the posthumous events that the wish concerns bear on less of what the person cared about, so by Principle~\ref{p:surv} their effect on how her life went is smaller.
\end{proof}

\begin{theorem}[Duties to the dead]\label{t:main}
We owe the dead fidelity to their reasonable wishes, and this duty weakens over time.
\end{theorem}

\begin{proof}
By Lemmas~\ref{l:subject} and~\ref{l:content}, we owe the dead, in the person of the ante-mortem person, fidelity to their reasonable wishes. By Principle~\ref{p:rel}, the obligation persists until discharged, released or defeated, and it is not released, since the dead cannot release. By Lemma~\ref{l:weak}, the interests it protects become fewer and weaker as time passes, so by Assumption~\ref{a:weight} its weight against competing reasons diminishes. The claims of the living, which compete with it, do not diminish. So the duty is defeated in more and more cases as time passes: it weakens.
\end{proof}

Theorem~\ref{t:main} explains a feature of the law that is otherwise puzzling: the rules that limit how long the wishes of a testator can govern property, and the doctrines that allow courts to apply charitable bequests to new purposes when their original purpose has become impossible, are not arbitrary limits on what the dead are owed, but recognition that what they are owed weakens as the interests behind it are spent.

\section{Objections and replies}

\begin{objection}[No subject]\label{o:nosub}
The dead do not exist, so nothing can be owed to them, and the appearance of posthumous duty is a projection of our concern for those we remember \citep{Callahan1987}.
\end{objection}

\begin{reply}
The duty is owed to someone who existed, the ante-mortem person (Lemma~\ref{l:subject}). The objection holds only if a person can be owed something only at times at which she exists, and that is the strong existence requirement that EXI-02 rejects. The projection account also misdescribes the cases: a deathbed promise made in secret to someone no living person remembers still binds, and nothing in the concern of the living explains why.
\end{reply}

\begin{objection}[Persistence]\label{o:persist}
ETH-16 holds that obligations grounded in reliance persist when the psychological states that the promise produced cease. The dead have no psychological states at all; on that view, nothing changes with the passage of time, and the duty should not weaken.
\end{objection}

\begin{reply}
Theorem~\ref{t:main} agrees that the duty persists: it is not released and it does not lapse when no one relies on it any longer, as Principle~\ref{p:rel} requires. What changes is its weight, and ETH-16 explicitly leaves the stringency of reliance-based obligations open. The weakening of Lemma~\ref{l:weak} is not a decline in anyone's state of relying; it is the fulfilment, extinction or obsolescence of the interests the entitlement protects, which is a change in what keeping the obligation would achieve for the person to whom it is owed.
\end{reply}

\begin{objection}[Only the living]\label{o:part}
The practice of honoring the wishes of the dead is justified by its value to the living, who want assurance that their own wishes will be honored. There is no need to posit duties to the dead \citep{Partridge1981}.
\end{objection}

\begin{reply}
Partridge identifies the right mechanism and the wrong beneficiary. The assurance the practice gives to the living creates their entitlement, and the entitlement is theirs; when they die it is not transferred to the community that gave it. If the practice were justified only by its value to those still living, a secret breach that could never be discovered would wrong no one, which is the conclusion Lemma~\ref{l:subject} avoids.
\end{reply}

\begin{objection}[Deathbed promises]\label{o:promise}
A deathbed promise to scatter someone's ashes at sea binds more strongly than a wish expressed in a will, and does not weaken.
\end{objection}

\begin{reply}
A deathbed promise adds to the general reliance of Assumption~\ref{a:practice} a specific reliance created by a particular person, and by Principle~\ref{p:rel} the promisor is obligated in his own right; the promise binds more. But its weight still follows the interests it protects. A promise to scatter ashes is usually kept within weeks, before any weakening; a promise to tend a grave for ever does weaken, as the interest in the grave's being tended becomes an interest in the regard of people who no longer exist.
\end{reply}

\section{Conclusion}

We owe the dead fidelity to their reasonable wishes. The duty is owed to the person as she was while she lived, in whom the practice of honoring such wishes created a reasonable reliance, and it protects her interests in what happens after her death. It persists, as reliance-based obligations do, but its weight follows the interests it protects, and these are fulfilled, extinguished or overtaken as time passes. The duty therefore weakens over time.

\begin{thebibliography}{99}\small
\bibitem[Aristotle(NE)]{AristotleNE} Aristotle. \emph{Nicomachean Ethics}.
\bibitem[Callahan(1987)]{Callahan1987} Callahan, J.~C. (1987). On harming the dead. \emph{Ethics}, 97(2), 341--352.
\bibitem[Feinberg(1984)]{Feinberg1984} Feinberg, J. (1984). \emph{Harm to Others}. Oxford University Press.
\bibitem[Partridge(1981)]{Partridge1981} Partridge, E. (1981). Posthumous interests and posthumous respect. \emph{Ethics}, 91(2), 243--264.
\bibitem[Pitcher(1984)]{Pitcher1984} Pitcher, G. (1984). The misfortunes of the dead. \emph{American Philosophical Quarterly}, 21(2), 183--188.
\end{thebibliography}

\end{document}
